\anexo{anexo:actions}{Acciones del modelo}

En un turno básico de batalla, el jugador puede escoger entre 14 acciones posibles.

\[
    a_t^i =
    \begin{cases}
        $\text{Movimiento 1}$,                      \\
        $\text{Movimiento 1 + Teracristalización}$, \\
        $\text{Movimiento 2}$,                      \\
        $\text{Movimiento 2 + Teracristalización}$, \\
        $\text{Movimiento 3}$,                      \\
        $\text{Movimiento 3 + Teracristalización}$, \\
        $\text{Movimiento 4}$,                      \\
        $\text{Movimiento 4 + Teracristalización}$, \\
        $\text{Cambio a Pokémon 0}$,                \\
        $\text{Cambio a Pokémon 1}$,                \\
        $\text{Cambio a Pokémon 2}$,                \\
        $\text{Cambio a Pokémon 3}$,                \\
        $\text{Cambio a Pokémon 4}$,                \\
        $\text{Cambio a Pokémon 5}$.                \\
    \end{cases}
\]

Entonces el vector de salida de la red neuronal va a tener una dimensión de 14. El entorno nos dice que acciones son legales y se enmascara con un $-\infty$ las que no lo son. Se aplica un softmax y se obtiene la probabilidad de cada acción. Se samplea una acción y se devuelve al entorno.

La teracristalización es la mecánica característica de la IX generación de Pokémon. Se puede hacer una vez por batalla y solo se puede hacer con un movimiento, no con un cambio. Por tanto, si el jugador ya ha teracristalizado, las acciones que incluyen la teracristalización no van a ser legales.

La librería poke-env~\cite{poke_env} siempre devuelve el equipo en el mismo orden, por lo que el Pokémon 0 es el primero que también aparece en los inputs de la red neuronal. Una restricción es que no se puede cambiar al Pokémon activo, por lo que si el Pokémon activo es el 0, la acción de cambio a Pokémon 0 no va a ser legal.
